\begin{figure}[H]
\centering
\begin{tikzpicture}[
  >=Latex,
  every node/.style={font=\small},
  elbo/.style={
    draw,
    rounded corners=2pt,
    align=center,
    inner sep=7pt,
    text width=12.2cm
  },
  view/.style={
    draw,
    rounded corners=2pt,
    align=center,
    inner sep=7pt,
    text width=5.75cm,
    minimum height=2.35cm
  },
  flow/.style={->,thick}
]

\node[elbo] (joint) at (0,2.25) {
  \textbf{联合密度形式}\\[4pt]
  $\displaystyle
  \mathcal{L}(\mathbf{x};\theta,\phi)
  =
  \mathbb{E}_{q_{\phi}(\mathbf{z}\mid\mathbf{x})}
  \!\left[
  \log
  \frac{p_{\theta}(\mathbf{x},\mathbf{z})}
       {q_{\phi}(\mathbf{z}\mid\mathbf{x})}
  \right]$
};

\node[view] (inf) at (-3.45,-1.2) {
  \textbf{推断视角}\\[4pt]
  $\displaystyle
  \begin{aligned}
  \mathcal{L}
  &=
  \log p_{\theta}(\mathbf{x})\\[-1pt]
  &\quad-
  \mathrm{KL}\!\left(
  q_{\phi}(\mathbf{z}\mid\mathbf{x})
  \,\|\,p_{\theta}(\mathbf{z}\mid\mathbf{x})
  \right)
  \end{aligned}$\\[5pt]
  精确对数似然与 ELBO 之间的差距来自
  近似后验与真实后验的不匹配。
};

\node[view] (train) at (3.45,-1.2) {
  \textbf{训练视角}\\[4pt]
  $\displaystyle
  \begin{aligned}
  \mathcal{L}
  &=
  \mathbb{E}_{q_{\phi}(\mathbf{z}\mid\mathbf{x})}
  [\log p_{\theta}(\mathbf{x}\mid\mathbf{z})]\\[-1pt]
  &\quad-
  \mathrm{KL}\!\left(
  q_{\phi}(\mathbf{z}\mid\mathbf{x})
  \,\|\,p(\mathbf{z})
  \right)
  \end{aligned}$\\[5pt]
  期望重构对数似然与
  朝先验的正则化相互权衡。
};

\draw[flow] (joint) -- node[left] {\scriptsize Bayes 规则} (inf);
\draw[flow] (joint) -- node[right] {\scriptsize 分解联合分布} (train);

\node[font=\normalsize\bfseries] at (0,3.45)
  {同一个 ELBO 的三种视角};

\end{tikzpicture}
\caption{\textbf{VAE ELBO 的三种等价视角。}
联合密度形式是一般变分目标。利用 Bayes 规则可得到推断视角，即“精确对数似然减去后验不匹配”；对联合分布进行分解则得到训练视角，即“期望重构对数似然减去朝潜变量先验的正则化”。三种形式在代数上完全等价。}
\label{fig:vae-elbo-three-views}
\end{figure}